% =============================================================
%  Section: Introduction
% =============================================================

Object detection is one of the canonical problems of computer vision. Given
an input image, the task is to enumerate the objects present, classify each
of them into one of a fixed set of categories, and localise each with a
bounding box. The output is therefore variable-length and structured, which
distinguishes detection from image classification (fixed-length output) and
from semantic segmentation (one prediction per pixel). The combination of
classification and localisation has been a productive testbed for
representation learning since the resurgence of deep convolutional networks
in the early 2010s~\citep{krizhevsky2012imagenet,russakovsky2015imagenet},
and remains a core capability underlying applications from autonomous
driving to medical imaging to industrial inspection.

% -------------------------------------------------------------
\subsection{Historical context}
\label{sec:intro:history}

Two algorithmic families have dominated the deep-learning era of object
detection, with markedly different trade-offs in accuracy, speed and
implementation complexity.

\paragraph{Two-stage detectors.}
The first wave, beginning with R-CNN~\citep{girshick2014rich}, decomposes the
problem into two sequential stages: first generate a set of candidate
regions where objects might be present, then classify each candidate
independently with a convolutional network. R-CNN's region proposal
mechanism (selective search) was external and expensive; its successors
Fast R-CNN~\citep{girshick2015fast} and
Faster R-CNN~\citep{ren2015faster} progressively absorbed the proposal step
into the network itself, trading the discrete proposal cascade for a learnt
region proposal network and shared convolutional features. Two-stage
detectors typically achieve the highest localisation accuracy on standard
benchmarks, at the cost of an inference path that touches every proposal at
least twice.

\paragraph{Single-stage detectors.} The second wave, of which
YOLOv1~\citep{redmon2016you} and SSD~\citep{liu2016ssd} are the canonical
representatives, reformulates detection as a single dense regression problem
on a fixed spatial grid. The network produces, in a single forward pass,
a complete tensor whose channels jointly encode class scores, box
parameters and per-box confidence at every cell of the grid. There is no
proposal stage and no per-region classifier; the entire predicted set of
detections is read off a single output tensor and filtered by non-maximum
suppression. Single-stage detectors are typically faster than two-stage
detectors by an order of magnitude and amenable to real-time deployment;
the original YOLO paper highlights this trade-off in its title.

The success of YOLO sparked a sustained research effort that produced
YOLOv2, YOLOv3~\citep{redmon2018yolov3} and a long line of community
descendants, each incorporating ideas (anchor boxes, multi-scale prediction,
fully convolutional heads, focal loss, mosaic augmentation) that
incrementally close the accuracy gap with two-stage detectors while
preserving the single-pass inference structure. By the time of writing, the
YOLO family remains the practical default for real-time detection on commodity
hardware.

% -------------------------------------------------------------
\subsection{Why study YOLOv1 specifically}
\label{sec:intro:why-yolov1}

This work targets the original YOLOv1 paper for three reasons.

First, it is the simplest detector in the single-stage family. The network
is a feedforward stack of convolutions followed by two dense layers; the
output is a single tensor; the loss is a single sum-of-squared-errors
expression with five clearly delineated terms. Everything that descendants
add (anchors, FPN, multi-scale outputs, focal loss) is absent. A successful
reproduction therefore exercises every concept that subsequent work assumes
the reader already understands, without yet introducing the optimisations
that obscure them.

Second, the paper is unusually clear about its own limitations. Section~4.4
of \citet{redmon2016you} enumerates failure modes (clusters of small objects,
unusual aspect ratios, coarse spatial localisation) that are inherent to
the YOLOv1 formulation and that are still useful to study in their original
form, because they motivate every architectural change in the family's later
versions.

Third, the implementation effort is small enough to be auditable
end-to-end by a single person, which is the prerequisite for the kind of
bug-forensics narrative that this work emphasises in
Section~\ref{sec:bug-investigation}. The audit identifies two implementation
mistakes that the unit-test suite did not catch, and uses them to argue for
a specific style of quantitative test that production deep-learning codebases
benefit from but tend not to write.

% -------------------------------------------------------------
\subsection{Objectives of this work}
\label{sec:intro:objectives}

The objectives are deliberately reproductive rather than novel. They are:

\begin{enumerate}
    \item \emph{Reproduce the YOLOv1 pipeline end-to-end in PyTorch.} The
        full system---dataset class, model, loss, distributed training loop
        with mixed precision, mAP evaluator and real-time inference
        script---is written from scratch and lives in the repository
        associated with this work. The implementation follows
        \citet{redmon2016you} on every point where the paper is explicit and
        is documented at the points where it is not.
    \item \emph{Port the paper's VOC setting ($C = 20$) to COCO 2017
        ($C = 80$).} The architectural changes are minimal (an output-layer
        widening and a per-cell tensor-layout adjustment) but the training
        regime is substantially harder, as Section~\ref{sec:disc:why-low}
        analyses in detail.
    \item \emph{Document the engineering reality of training a detector from
        scratch.} This includes the four-version iteration on hyper-parameters
        that did not solve the problem
        (Sections~\ref{sec:experiments:v1}--\ref{sec:experiments:v4}), the
        bug audit that did (Section~\ref{sec:bug-investigation}), the
        post-fix runs that produced the first non-zero metric
        (Sections~\ref{sec:experiments:v5}--\ref{sec:experiments:v6}), and
        the per-class breakdown that motivates the future-work programme of
        Section~\ref{sec:disc:future}.
\end{enumerate}

We do not claim competitiveness with current detectors, and we do not
present any methodological novelty. The contribution of this work is the
auditability of a complete YOLOv1 pipeline targeted at COCO, together with
the documentation of two failure modes that are easy to introduce, hard to
detect from training curves, and worth flagging to anyone reproducing a
detector for the first time.

% -------------------------------------------------------------
\subsection{Outline}
\label{sec:intro:outline}

Section~\ref{sec:theoretical_framework} formalises the problem and the
YOLOv1 loss, summarising the parts of \citet{redmon2016you} that subsequent
sections rely on. Section~\ref{sec:impl:dataset} describes the dataset
preparation pipeline; Section~\ref{sec:impl:model} introduces the model
architecture as a single TikZ diagram; Sections~\ref{sec:impl:loss}
through~\ref{sec:impl:eval} cover the loss implementation, the training
infrastructure and the evaluation pipeline. Section~\ref{sec:experimentation}
is the substantive narrative of the work: six training runs, the failure
of the first four, the audit that explained the failures, and the two
runs that produced the first non-zero detection metric. Section~\ref{sec:results}
reports the headline figures, the per-class breakdown, inference latency and
a side-by-side comparison with the original paper. Section~\ref{sec:disc:why-low}
discusses the limitations and the natural continuation of the work, and the
Conclusions restate the central finding. Two appendices contain the full
per-class AP table and a side-by-side listing of the loss function before
and after the index fix.
